\documentclass[12pt,aspectratio=1610,professionalfont,serif,t]{beamer}
\usepackage[osf,p]{newpxtext}
\usepackage{newpxmath}
\usepackage{fontawesome5}

\usepackage{fancyvrb}
\fvset{%frame=single,
commandchars=\\\{\},
%framesep=1mm,
fontfamily=helvetica,
fontsize=\normalsize
}
%\usepackage[bitstream-charter]{mathdesign}
\usepackage{listings}
\lstset{commentstyle=\color{orange}, keywordstyle=\color{yellow} }


\usetheme{Boadilla}
\useoutertheme{split}
\usecolortheme{albatross}


\setbeamertemplate{blocks}[rounded][shadow=false]
\setbeamersize{text margin left=\dimexpr 1cm+1ex\relax, text margin right=\dimexpr 1cm+1ex\relax}
\addtobeamertemplate{headline}{}{\usebeamerfont{frametitle}\vspace*{1ex}}
\setbeamertemplate{navigation symbols}{}
\setbeamertemplate{page number in head/foot}{}
%\beamerdefaultoverlayspecification{<+->}


\setbeamercolor*{structure}{fg=green!60!black,bg=black}
\setbeamercolor*{normal text}{fg=green!60!black,bg=black}
\setbeamercolor{palette primary}{use={structure,normal text},fg=green,bg=structure.bg}
\setbeamercolor{palette secondary}{use={structure,normal text},fg=structure.fg,bg=structure.bg}
\setbeamercolor{palette tertiary}{use={structure,normal text},fg=structure.fg,bg=structure.bg}
\setbeamercolor{palette quaternary}{use={structure,normal text},fg=white,bg=green!30!black}
\setbeamercolor{palette sidebar quaternary}{use={structure,normal text},fg=white,bg=green!30!black}
\setbeamercolor{title in head/foot}{use={structure,normal text},fg=white,bg=green!30!black}

\setbeamercolor*{example text}{use={structure},fg=green!60!black}
\setbeamercolor*{block body}{use={structure},bg=structure.bg}
\setbeamercolor*{block body alerted}{use={structure},bg=structure.bg}
\setbeamercolor*{block body example}{use={structure},bg=structure.bg}
\setbeamercolor*{block title}{parent=structure,bg=structure.bg}
\setbeamercolor*{block title alerted}{use={structure,alerted text},fg=alerted text.fg,bg=structure.bg}
\setbeamertemplate{navigation symbols}{}
\setbeamertemplate{items}[square]
\setbeamercolor{item projected}{fg=white}
\setbeamercolor*{normal text}{use={structure},fg=white!90!blue,bg=structure.bg}
\setbeamercolor*{separation line}{}
\setbeamercolor*{fine separation line}{}
\setbeamercolor{alerted text}{use={structure},fg=structure.fg}

\usepackage[italian]{babel}
\usepackage[utf8]{inputenc}
\usepackage{pgf}
\usepackage{verbatim}
\usepackage{inconsolata}
\usepackage{listings}
\lstset{language=C, frame=single,
	basicstyle=\ttfamily,
	numbers=left, numberstyle=\tiny\color{gray},
	numbersep=5pt, fancyvrb=true
}

\usepackage[noend]{algorithm2e}
\usepackage{pgf}
\usepackage{tikz}
\usepackage{graphicx}
\usetikzlibrary{%
	arrows,
	arrows.meta,
	positioning,
	calc,
	backgrounds,
	chains,
	matrix,
	patterns,
	automata,
	fit,
	graphs,
	decorations,
	decorations.pathmorphing,
	decorations.pathreplacing,
	decorations.markings,
}

\usepgflibrary{shapes,shapes.geometric}


\usepackage{url}


\usepackage{xmpmulti}
\usepackage{transparent}
\graphicspath{{img/}}

\newcommand{\zero}{\ensuremath{\mathtt{0}}\xspace}
\newcommand{\uno}{\ensuremath{\mathtt{1}}\xspace}
\newcommand{\due}{\ensuremath{\mathtt{2}}\xspace}
\newcommand{\tre}{\ensuremath{\mathtt{3}}\xspace}
\newcommand{\kleene}{\ensuremath{{}^{\ast}}\xspace}
\xspaceaddexceptions{\zero \uno \due \tre \kleene}


\newcommand{\simplebox}[1]{
      \begin{tikzpicture}
              \node[text centered, rectangle, draw=black, thick, text width=\linewidth,align=center, rounded corners, minimum height=1em] (C) {#1};
      \end{tikzpicture}
}
%
% In questo file si mettono i comandi utili solo per questo insegnamento.
%
\author{Gianluca Della Vedova}
\title{Linguaggi e Computabilità}
\institute{Univ. Milano--Bicocca\\
	\texttt{https://www.unimib.it/gianluca-della-vedova}}
\date{\today}



\begin{document}


\begin{frame}[fragile]
        \frametitle{Linguaggi e Computabilità: Lezione 18}
\begin{enumerate}
\item
Equivalenza dei PDA che accettano per pila vuota o per stato accettante
\end{enumerate}
\end{frame}


\begin{frame}[fragile]
        \frametitle{Costruire il DFA equivalente}
\begin{tikzpicture}[
    shorten >=1pt,
    node distance=2cm,
    on grid,
    auto,
    every state/.style={draw=black, thick, fill=white, minimum size=1.2cm},
    >=Stealth,
    thick
]

   % Define the states based on the layout
   \node[state, accepting, initial, initial text=] (q0) {$q_0$};
   \node[state] (q1) [below=of q0] {$q_1$};
   \node[state] (q3) [right=of q0, xshift=3cm] {$q_3$};
   \node[state, accepting] (q2) [below=of q3] {$q_2$};
   \node[state] (q4) [below left=of q1] {$q_4$};
   \node[state] (q5) [below right=1cm and 2.7cm of q3] {$q_5$};

   % Define the transition paths
   \path[->]
    % Transitions from q0
    (q0) edge [bend left=15] node {\zero,\uno} (q3)
    (q0) edge [bend left=20] node [right] {$\epsilon$} (q1)

    % Transitions from q1
    (q1) edge [bend left=20] node [left] {\zero} (q0)
    (q1) edge node [above] {$\varepsilon$} (q4)
    (q1) edge node [above left] {\uno} (q2)

    % Transitions from q2
    (q2) edge [out=165, in=-15, looseness=.4] node [above left, yshift=1ex] {\zero} (q0)
    (q2) edge [bend right=20] node [right] {\uno} (q3)

    % Transitions from q3
    (q3) edge [bend right=20] node [left] {\zero} (q2)

    % Transitions from q4
    (q4) edge [out=75, in=195, looseness=0.9] node [left] {\zero} (q0)

    % Transitions from q5
    (q5) edge node [below right] {\zero} (q2)
    (q5) edge node [above right] {\uno} (q3)
    (q5) edge [out=270, in=15] node [above] {\zero} (q4);

\end{tikzpicture}
\end{frame}

\begin{frame}[fragile]
        \frametitle{Pumping Lemma}
        $L = \{ \uno^{2n}\zero\uno\zero\zero^{n-1}\uno^{n} \; \vert \; n \ge 1 \}$ è regolare ?
        \end{frame}

\begin{frame}[fragile]
\frametitle{Costruire il DFA}
linguaggio $L$ sull'alfabeto $\Sigma = \{\zero, \uno, \due, \tre\}$ e formato dalle
parole $w$ tali che:
(1) $w$ contiene la sottostringa \zero\zero\uno;
(2) $w$ termina con \zero;
(3) $w$ inizia con almeno due \zero.\\
Ad esempio, le parole \zero\zero\uno\zero, \zero\zero\uno\zero\uno\zero,
\zero\zero\uno\uno\due\uno\zero,
\zero\zero\zero\tre\zero\zero\uno\due\tre\zero appartengono a $L$,
mentre
le parole \zero\zero\due\uno\due\tre\zero, \zero\zero\uno\due\uno\tre\tre, \zero\uno\due\uno\tre\uno\zero non appartengono a $L$.
\end{frame}

\end{document}
